% % !TeX root = Notes.tex

% \documentclass{article}
% \usepackage{ctex}
% \usepackage[utf8]{inputenc}
% \usepackage[top=1in, bottom=1in, left=1.in, right=1.in]{geometry} % adjust default margins
% \usepackage{fancyhdr}
% \pagestyle{fancy}
% \lhead{Yiran Wang}\chead{\reportNumber}\rhead{Rutgers Univ.}\lfoot{}\cfoot{\thepage}\rfoot{} % setting page header and footer
% \setlength{\parindent}{0em}\setlength{\parskip}{0.5ex} % use newline to seperate paragraphs
% \usepackage[T1]{fontenc}     % oriented to output, that is, what fonts to use for printing characters % [must use] make ligatures copyable
% \usepackage{enumitem} %\setlist[itemize]{noitemsep,nolistsep} % same to use \begin{itemize}[nolistsep,noitemsep]
% \usepackage[ampersand]{easylist} % ampersand == & % \begin{easylist}[itemize] %% to make easy list environment
% \usepackage{amsmath,amssymb,bm}
% \usepackage[colorlinks=true,linkcolor=blue,citecolor=blue,urlcolor=blue,breaklinks=true]{hyperref} % color the hyperlink rather than box it
% \usepackage{graphicx}
% \usepackage{color}
% \usepackage[all]{hypcap} % make the hyperlink to point at the upper border of a figure of table, rather than to point to its caption.
% \usepackage{breakcites}
% \usepackage{cleveref} \crefname{equation}{Eq.}{Eqs.} \Crefname{equation}{Equation}{Equations} \crefname{figure}{Fig.}{Figs.} \Crefname{figure}{Figure}{Figures} \crefname{table}{Table}{Tables} % allow refer to multiple equations and change cross-reference name
% \usepackage{caption}
% \usepackage{subfigure}
% \usepackage{wasysym} % to insert emotion \smiley
% \usepackage{float} 
% \usepackage[yyyymmdd,hhmmss]{datetime}\renewcommand{\dateseparator}{-}
% \usepackage{placeins}
% \usepackage{setspace} 
% \singlespacing 
% % \setstretch{0.9}
% % \doublespacing \onehalfspacing % 调整行距
% %\usepackage{pdflscape} \horizontal figure
% \usepackage{ulem}
% \usepackage{mathtools}
% \urlstyle{same}
% %\graphicspath{ {./Test_01/} }
% \bibliographystyle{abbrv}
% %%%%%%%%%%%%%%%%%%%%%%%%%%%%%%%%%%%%%%%%%%%%%%%%%%%%%%%
% \newcommand{\Yiran}[1]{\textcolor{blue}{#1}}
% %%%%%%%%%%%%%%%%%%%%%%%%%%%%%%%%%%%%%%%%%%%%%%%%%%%%%%%

% \title{Ads Notes}
% \author{\mathfrak{Y} \mathfrak{R} }
% \date{Updated on \today}

% !TeX root = Notes.tex

\documentclass{article}

%%%%%%%%%%%%%%%%%%%%%%%%%%%%%%%%%%%%%%%%%%%%%%%%%%%%%%%
% Language
%%%%%%%%%%%%%%%%%%%%%%%%%%%%%%%%%%%%%%%%%%%%%%%%%%%%%%%
\usepackage{ctex}
\usepackage[utf8]{inputenc}
\usepackage[T1]{fontenc}

%%%%%%%%%%%%%%%%%%%%%%%%%%%%%%%%%%%%%%%%%%%%%%%%%%%%%%%
% Page Layout
%%%%%%%%%%%%%%%%%%%%%%%%%%%%%%%%%%%%%%%%%%%%%%%%%%%%%%%
\usepackage[
    top=1in,
    bottom=1in,
    left=1in,
    right=1in
]{geometry}

%%%%%%%%%%%%%%%%%%%%%%%%%%%%%%%%%%%%%%%%%%%%%%%%%%%%%%%
% Header & Footer
%%%%%%%%%%%%%%%%%%%%%%%%%%%%%%%%%%%%%%%%%%%%%%%%%%%%%%%
\usepackage{fancyhdr}
\pagestyle{fancy}
\lhead{Yiran Wang}
\chead{\reportNumber}
\rhead{Rutgers Univ.}
\lfoot{}
\cfoot{\thepage}
\rfoot{}

%%%%%%%%%%%%%%%%%%%%%%%%%%%%%%%%%%%%%%%%%%%%%%%%%%%%%%%
% Paragraph & Line Spacing
%%%%%%%%%%%%%%%%%%%%%%%%%%%%%%%%%%%%%%%%%%%%%%%%%%%%%%%
\usepackage{setspace}
\setstretch{0.95}              % slightly tighter than default

\setlength{\parindent}{0em}
\setlength{\parskip}{0.4ex}

%%%%%%%%%%%%%%%%%%%%%%%%%%%%%%%%%%%%%%%%%%%%%%%%%%%%%%%
% Lists
%%%%%%%%%%%%%%%%%%%%%%%%%%%%%%%%%%%%%%%%%%%%%%%%%%%%%%%
\usepackage{enumitem}
\setlist{
    noitemsep,
    topsep=2pt,
    parsep=0pt,
    partopsep=0pt
}

\usepackage[ampersand]{easylist}

%%%%%%%%%%%%%%%%%%%%%%%%%%%%%%%%%%%%%%%%%%%%%%%%%%%%%%%
% Math
%%%%%%%%%%%%%%%%%%%%%%%%%%%%%%%%%%%%%%%%%%%%%%%%%%%%%%%
\usepackage{amsmath,amssymb,bm}
\usepackage{mathtools}

% Reduce display math spacing
\setlength{\abovedisplayskip}{6pt}
\setlength{\belowdisplayskip}{6pt}
\setlength{\abovedisplayshortskip}{3pt}
\setlength{\belowdisplayshortskip}{3pt}

%%%%%%%%%%%%%%%%%%%%%%%%%%%%%%%%%%%%%%%%%%%%%%%%%%%%%%%
% Figures
%%%%%%%%%%%%%%%%%%%%%%%%%%%%%%%%%%%%%%%%%%%%%%%%%%%%%%%
\usepackage{graphicx}
\usepackage{subfigure}
\usepackage{caption}
\usepackage{float}
\usepackage{placeins}

%%%%%%%%%%%%%%%%%%%%%%%%%%%%%%%%%%%%%%%%%%%%%%%%%%%%%%%
% Hyperlinks
%%%%%%%%%%%%%%%%%%%%%%%%%%%%%%%%%%%%%%%%%%%%%%%%%%%%%%%
\usepackage[
    colorlinks=true,
    linkcolor=blue,
    citecolor=blue,
    urlcolor=blue,
    breaklinks=true
]{hyperref}

\usepackage[all]{hypcap}
\usepackage{breakcites}
\urlstyle{same}

%%%%%%%%%%%%%%%%%%%%%%%%%%%%%%%%%%%%%%%%%%%%%%%%%%%%%%%
% References
%%%%%%%%%%%%%%%%%%%%%%%%%%%%%%%%%%%%%%%%%%%%%%%%%%%%%%%
\usepackage{cleveref}

\crefname{equation}{Eq.}{Eqs.}
\Crefname{equation}{Equation}{Equations}

\crefname{figure}{Fig.}{Figs.}
\Crefname{figure}{Figure}{Figures}

\crefname{table}{Table}{Tables}

%%%%%%%%%%%%%%%%%%%%%%%%%%%%%%%%%%%%%%%%%%%%%%%%%%%%%%%
% Misc
%%%%%%%%%%%%%%%%%%%%%%%%%%%%%%%%%%%%%%%%%%%%%%%%%%%%%%%
\usepackage{color}
\usepackage{wasysym}
\usepackage[yyyymmdd,hhmmss]{datetime}
\renewcommand{\dateseparator}{-}
\usepackage{ulem}


%%%%%%%%%%%%%%%%%%%%%%%%%%%%%%%%%%%%%%%%%%%%%%%%%%%%%%%
% Bibliography
%%%%%%%%%%%%%%%%%%%%%%%%%%%%%%%%%%%%%%%%%%%%%%%%%%%%%%%
\bibliographystyle{abbrv}

%%%%%%%%%%%%%%%%%%%%%%%%%%%%%%%%%%%%%%%%%%%%%%%%%%%%%%%
% Custom Commands
%%%%%%%%%%%%%%%%%%%%%%%%%%%%%%%%%%%%%%%%%%%%%%%%%%%%%%%
\newcommand{\Yiran}[1]{\textcolor{blue}{#1}}


\title{Ads Notes}
\author{YR}
\date{Updated on \today}


\begin{document}
\maketitle
\section{Ranking Score}

In ads ranking, the ranking score is often estimated as:
\[
\text{Ranking Score} = \text{Bid} \times pCTR \times pCVR
\]
which represents the expected value (e.g., expected revenue) per impression.

Auction: rank ads by score and determine which ads are shown and sometimes their prices. (展示广告的排序和位置)

User experiences will affect  when considering pCTR. 

(User experiences 会影响CTR, 但另一方面来说，不能说pCTR直接受user experience影响，因为“图片吸引人”或者“标题具有引导性”都会提成pCTR但是用户体验是negative的)


\section{Commonly used feature engineering}
\begin{itemize}
    \item user features
        \subitem historical CTR, demographics, user embedding
    \item ad features
        \subitem creative(image, text, video features), ad history, ad category , historical performance
    
    \item Historical statistical features
        \subitem user CTR over 1/7/30 days.
        \subitem ad CTR by category , user-specific ad CTR
    
    \item cross features
        \subitem user clicked the ads before? ( user $\times$ ad /category /advitiser)
        \subitem user-ad interaction frequency
    \item context features
        \subitem time, device, position, hour of the day, connection type
    \item features crossing for nonlinear relationship
        \subitem user age $\times$ ad category 
\end{itemize}

    
\section{常用训练模型}
Increasing capability to model feature interaction
\begin{itemize}
    \item Logistic Regression: Linear model, suitable for sparse features. Relay heavily on manual feature engineering
    \item LR+ feature crossing
    \item GBDT (Gradient Boosted Decision Trees): 优势： automatic fetaure partitioning, nonlinear interactions, no feature scaling.\\
    常用： XGboost, LightBGM
    \item GBDT + LR

    \item Deep CTR models: DeepFM / DCN (deep cross network). for low and high order interactions
    \item DIN (deep Interest Network): Models users interests from \textbf{historical behaviour sequences} using attention.\\
    eg: 点击历史，浏览历史，购买历史。
    \item DIEN/Transformer-based models
\end{itemize}

\section{Cold Start}
Modern approach: Content-based embeddings + representation learning

For new ads:
\begin{itemize}
    \item content based features
    \begin{itemize}
        \item categorical features: advertiser-ID, category , campaign type
        \item creative features: image(handle with CNN or Version Transformer), text embedding (ad title or description)
        \item format features: video/ image, color, placement(广告位),...
    \end{itemize}

    
    \item leverage similar ads
    \begin{itemize}
        \item build ad embedding share the same or similar category /advertiser-ID
        \item use KNN in embedding space
    \end{itemize}

    \item exploration strategy
    \begin{itemize}
        \item user Thompson Sampling or Upper Confidence Bound (UCB) to initially show the ad to small and representative sample
        \item start with prior  CTR (like category average value) and update 
        \item For first 1000 impression, for example, on "exploration", then on real data.
    \end{itemize}
\end{itemize}

For new user
\begin{itemize}
    \item contextual signals and request time
    \begin{itemize}
        \item Device models, OS, network type. time of day, location, languages, traffic source( from website/app/events...)
    \end{itemize}

    \item lookalike model
    \begin{itemize}
        \item nearest cluster of existing user
    \end{itemize}

    \item rapid adaptation: 新用户虽然一开始没有历史数据，但只要产生少量行为，就应该尽快利用这些行为更新模型对该用户的理解，而不是一直依赖默认画像。
    \begin{itemize}
        \item Continuously update user embeddings as interactions arrive.
    \end{itemize}
\end{itemize}

Embeddings can be learned from historical click data using two-tower models, Word2Vec-style objectives, or deep CTR models.

\section{案例分析： AUC improved  but revenue dropped}
Suppose AUC improved 2\% but revenue droppoed 1\%. What will you do.

\begin{itemize}
    \item Immediate check:
    \begin{itemize}
        \item Data leaks?
        \item verify feature distribution
        \item check feature freshments 
    \end{itemize}

    \item Model behaviour analysis
    \begin{itemize}
        \item calibration check
        \item segment-level analysis:
        \begin{itemize}
            \item Ad catagory: over predicted cheap ads
            \item User segment: how traffic allocated to high/low vaue users
            \item traffic sources: IOS/Android, different country. Compared with trained data
        \end{itemize}
        \item position bias
    \end{itemize} 

    \item Auction Dynamics
    \begin{itemize}
        \item Model may over-predicted CTR for low-bid ads, and under-predicted CTR for high-bid ads.（训练时便宜的多，贵的少，导致实际收益下降）
        \item If in section-price auction, ranking = pCTR * bid. 
        \item 
    \end{itemize}
\end{itemize}
% \bibliography{reference}
\end{document}